\clearpage
\VerdierPhysicalPageAnchor{0170}%
\begingroup\small\setlength{\parskip}{0pt}\setlength{\emergencystretch}{2.5em}
\VerdierSemanticAnchor{ch4-4-2-7-second-diagram}
En appliquant la même remarque au diagramme :
\[
\begin{tikzcd}[column sep=4.6em,row sep=2.8em]
&H^n(f_2f_1)\arrow[d]\arrow[dr]&\\
H^n(f_1)\arrow[ur]\arrow[r]&
H^n(f_3f_2f_1)\arrow[r]&H^n(f_3f_2),
\end{tikzcd}
\]
on obtient un isomorphisme :
\VerdierSemanticAnchor{ch4-4-2-7-2}
\[
E\!\begin{bmatrix}n\\ f_1,f_2,f_3\end{bmatrix}
\xrightarrow{\ \sim\ }
\operatorname{Im}\bigl(H^n(f_2f_1)\longrightarrow H^n(f_3f_2f_1)\bigr)
\Big/
\operatorname{Im}\bigl(H^n(f_1)\longrightarrow H^n(f_3f_2f_1)\bigr).
\tag{4.2.7.2}
\]

\VerdierSemanticAnchor{ch4-4-3}
\medskip\noindent\textbf{4.3. Les termes
$E\!\begin{bmatrix}n\\ -\infty,p,q,+\infty\end{bmatrix}$.}

\VerdierSemanticAnchor{ch4-4-3-1}
\medskip\noindent\textbf{4.3.1.}
Soit $\widetilde{\mathbb Z}$ l'ensemble ordonné $\mathbb Z$ complété par un
élément initial $-\infty$ et un élément final $+\infty$. Par abus de
notation, nous désignerons également par $\widetilde{\mathbb Z}$ la
catégorie associée à l'ensemble ordonné $\widetilde{\mathbb Z}$. Un
morphisme de $\widetilde{\mathbb Z}$ est donc un couple $p\leq q$, où $p$
et $q$ sont des éléments de $\widetilde{\mathbb Z}$. Une suite ordonnée
$p_1\leq p_2\leq\cdots\leq p_n$ définit un objet de
$\mathcal{F}\ell_{n-1}(\widetilde{\mathbb Z})$ noté
$(p_1,p_2,\ldots,p_n)$. Un morphisme de $\widetilde{\mathbb Z}$ est dit
\textit{fini} si sa source et son but appartiennent à $\mathbb Z$. Un objet
de $\mathcal{F}\ell_n(\widetilde{\mathbb Z})$ est dit \textit{fini} si les
morphismes qui le composent sont finis.

\VerdierSemanticAnchor{ch4-4-3-2}
\medskip\noindent\textbf{4.3.2.}
\textit{Nous n'utiliserons dans ce travail que les objets spectraux de type
$\widetilde{\mathbb Z}$, qui seront appelés dorénavant objets spectraux.}
On retrouve ainsi la notion du chapitre XV, §7 de [1]. Soit
$(H^n(p,q),\delta,\ p\leq q\in\widetilde{\mathbb Z})$ un objet spectral à
valeurs dans une catégorie abélienne $\mathcal A$. Posons :
\VerdierSemanticAnchor{ch4-4-3-2-1}
\[
F^pH^n(-\infty,+\infty)
=\operatorname{Im}\bigl(H^n(-\infty,p)\longrightarrow
H^n(-\infty,+\infty)\bigr).
\tag{4.3.2.1}
\]
On a $F^{+\infty}H^n(-\infty,+\infty)=H^n(-\infty,+\infty)$ et
$F^{-\infty}H^n(-\infty,+\infty)=0$. Les
$F^pH^n(-\infty,+\infty)$ forment une filtration croissante de type
$\widetilde{\mathbb Z}$ de $H^n(-\infty,+\infty)$. De plus, pour tout
couple $p\leq q$ on a un isomorphisme (4.2.7.2) :
\VerdierSemanticAnchor{ch4-4-3-2-filtration-quotient}
\[
E\!\begin{bmatrix}n\\ -\infty,p,q,+\infty\end{bmatrix}
\simeq
F^qH^n(-\infty,+\infty)\big/F^pH^n(-\infty,+\infty).
\]
\endgroup
